\documentclass[12pt,a4paper]{article}

\usepackage[utf8]{inputenc}
\usepackage[T1]{fontenc}
\usepackage{amsmath}
\usepackage{amssymb}
\usepackage{booktabs}
\usepackage{array}
\usepackage{geometry}
\usepackage{hyperref}
\usepackage{cite}

\geometry{
  a4paper,
  top=2.5cm,
  bottom=2.5cm,
  left=2.5cm,
  right=2.5cm
}

\hypersetup{
  colorlinks=true,
  linkcolor=blue,
  citecolor=blue,
  urlcolor=blue
}

\title{\textbf{Deuterium-Mediated Neutron Chain Fusion Initiated\\
by Proton--Electron Capture with Lithium-6 Yield Enhancement:\\
A Tritium Breeding Mechanism for Existing Tokamak Infrastructure}}

\author{Eymen Aydoğan}

\date{May 2026}

\begin{document}

\maketitle

\begin{abstract}
We propose a theoretical mechanism in which free protons injected into
dense deuterium gas undergo Coulomb-barrier-free electron capture
($p + e^- \rightarrow n + \nu_e$), producing neutrons that
subsequently convert deuterium nuclei into tritium via
$n + {}^2\text{H} \rightarrow {}^3\text{H} + p + 2.22\,\text{MeV}$.
The liberated proton re-enters the chain, sustaining neutron
production. Accumulated tritium and deuterium then undergo
fusion (${}^3\text{H} + {}^2\text{H} \rightarrow {}^4\text{He}
+ n + 17.6\,\text{MeV}$), with Lithium-6 acting as a yield
enhancer by capturing residual neutrons
($n + {}^6\text{Li} \rightarrow {}^4\text{He} + {}^3\text{H}
+ 4.78\,\text{MeV}$) and recycling tritium back into the fusion
cycle. The central contribution of this work is that the
proton--electron capture chain provides a self-sustaining
\emph{in situ} tritium breeding mechanism fully compatible with
existing tokamak infrastructure, addressing the critical tritium
supply problem in D--T fusion programs.
\end{abstract}

\noindent\textbf{Keywords:} tritium breeding, deuterium, proton-electron
capture, neutron chain, tokamak, ITER, lithium-6, D--T fusion

\hrule
\vspace{1em}

%----------------------------------------------------------------------
\section{Introduction}
%----------------------------------------------------------------------

Deuterium--tritium (D--T) fusion is the primary reaction pursued by
major tokamak programs worldwide, including ITER and China's EAST
reactor~\cite{iter2023,east2023}. The reaction:

\begin{equation}
  {}^2\text{H} + {}^3\text{H} \rightarrow {}^4\text{He} + n
  + 17.6\,\text{MeV}
  \label{eq:dtfusion}
\end{equation}

releases 17.6 MeV per event and is the most accessible fusion
reaction at achievable plasma temperatures ($T \sim 10^8$ K).
However, tritium (${}^3\text{H}$) does not occur naturally in
significant quantities and has a half-life of only 12.3 years,
making tritium supply one of the central unsolved problems of
fusion energy~\cite{abdou2021}.

Current tritium breeding concepts rely on neutron bombardment of
lithium blankets surrounding the reactor. While viable in principle,
this approach requires careful neutron economy management and does
not achieve tritium self-sufficiency without a breeding ratio
exceeding unity~\cite{abdou2021}.

In this paper, we propose a complementary mechanism: a
proton--electron capture chain in deuterium gas that produces
neutrons without a Coulomb barrier, converts deuterium to tritium
\emph{in situ}, and feeds the resulting tritium directly into the
D--T fusion cycle. Lithium-6 serves as a yield enhancer rather than
the primary energy source, capturing excess neutrons and regenerating
tritium.

%----------------------------------------------------------------------
\section{Reaction Mechanism}
%----------------------------------------------------------------------

\subsection{Step 1: Coulomb-Barrier-Free Neutron Production}

A free proton with kinetic energy $E_k \geq 1.293$ MeV is injected
into dense deuterium gas. The proton is attracted to the bound
electron of a deuterium atom (opposite charges, no Coulomb barrier):

\begin{equation}
  p + e^- \rightarrow n + \nu_e
  \qquad \Delta mc^2 = -1.293\,\text{MeV}
  \label{eq:capture}
\end{equation}

This is identical to the mechanism proposed in our companion
paper~\cite{aydogan2026a}, extended here to a deuterium medium.

\subsection{Step 2: Neutron--Deuterium Tritium Breeding}

The produced neutron interacts with a deuterium nucleus:

\begin{equation}
  n + {}^2\text{H} \rightarrow {}^3\text{H} + p + 2.22\,\text{MeV}
  \label{eq:breeding}
\end{equation}

This reaction has \textbf{no Coulomb barrier} (neutron is neutral),
proceeds with thermal neutrons, and has a well-measured cross-section
of $\sigma \approx 0.52$ mb at thermal energies~\cite{pdg2022}.
The liberated proton ($p$) re-enters the chain at Step~1,
sustaining neutron production without external proton injection
after initiation.

\subsection{Step 3: D--T Fusion}

Accumulated tritium undergoes fusion with deuterium:

\begin{equation}
  {}^3\text{H} + {}^2\text{H} \rightarrow {}^4\text{He} + n
  + 17.6\,\text{MeV}
  \label{eq:fusion}
\end{equation}

The Coulomb barrier for this step is overcome by the same methods
employed in existing tokamak reactors: magnetic confinement and
plasma heating to $T \sim 10^8$ K. \textbf{No new technology is
required for this step.} The novelty of this proposal lies in
providing the tritium fuel via Steps~1--2.

\subsection{Step 4: Lithium-6 Yield Enhancement}

Neutrons produced in Eq.~\eqref{eq:fusion} are captured by a
Lithium-6 blanket:

\begin{equation}
  n + {}^6\text{Li} \rightarrow {}^4\text{He} + {}^3\text{H}
  + 4.78\,\text{MeV}
  \label{eq:li6}
\end{equation}

The tritium produced here is recycled back into Step~3, increasing
the tritium breeding ratio. Lithium-6 is thus a \emph{yield
enhancer}, not the primary energy source.

%----------------------------------------------------------------------
\section{Complete Reaction Cycle}
%----------------------------------------------------------------------

The full cycle is:

\begin{align}
  p + e^- &\rightarrow n + \nu_e
  \quad -1.293\,\text{MeV} \tag{S1}\\
  n + {}^2\text{H} &\rightarrow {}^3\text{H} + p
  \quad +2.22\,\text{MeV} \tag{S2}\\
  {}^3\text{H} + {}^2\text{H} &\rightarrow {}^4\text{He} + n
  \quad +17.6\,\text{MeV} \tag{S3}\\
  n + {}^6\text{Li} &\rightarrow {}^4\text{He} + {}^3\text{H}
  \quad +4.78\,\text{MeV} \tag{S4}
\end{align}

Step S2 regenerates the proton for S1, closing the neutron
production loop. Step S4 regenerates tritium for S3, closing
the fusion fuel loop.

\subsection{Net Energy per Full Cycle}

\begin{table}[h!]
\centering
\caption{Energy budget per complete reaction cycle}
\vspace{0.5em}
\begin{tabular}{lcc}
\toprule
\textbf{Process} & \textbf{Energy} & \textbf{Direction} \\
\midrule
$p + e^- \rightarrow n + \nu_e$                          & 1.293 MeV  & Input  \\
$n + {}^2\text{H} \rightarrow {}^3\text{H} + p$         & 2.220 MeV  & Output \\
${}^3\text{H} + {}^2\text{H} \rightarrow {}^4\text{He} + n$ & 17.600 MeV & Output \\
$n + {}^6\text{Li} \rightarrow {}^4\text{He} + {}^3\text{H}$ & 4.780 MeV  & Output \\
\midrule
\textbf{Net per cycle} & \textbf{23.307 MeV} & \textbf{Output} \\
\bottomrule
\end{tabular}
\end{table}

\begin{equation}
  \boxed{E_\text{net} = 2.220 + 17.600 + 4.780 - 1.293
  = 23.307\,\text{MeV per cycle}}
\end{equation}

%----------------------------------------------------------------------
\section{Comparison with Existing Approaches}
%----------------------------------------------------------------------

\begin{table}[h!]
\centering
\caption{Comparison with existing tritium breeding and fusion approaches}
\vspace{0.5em}
\begin{tabular}{lccc}
\toprule
\textbf{Feature} & \textbf{ITER Blanket} & \textbf{Standard D--T} & \textbf{This Work} \\
\midrule
Tritium source       & Li blanket     & External supply  & \textbf{In situ chain} \\
Coulomb barrier      & N/A            & Yes (plasma)     & None (S1, S2)    \\
Neutron source       & Fusion         & Fusion           & p+e⁻ chain       \\
Li-6 role            & Primary breed  & None             & Yield enhancer   \\
Self-sustaining      & Partial        & No               & Yes (loop)       \\
Net energy (MeV)     & 17.6           & 17.6             & 23.307           \\
New infrastructure   & Required       & Required         & \textbf{None needed} \\
\bottomrule
\end{tabular}
\end{table}

The key advantage over ITER's blanket breeding approach is that
the proton--electron capture chain produces tritium \emph{before}
fusion occurs, rather than relying on fusion neutrons to breed
tritium after the fact. This eliminates the startup tritium
inventory problem.

%----------------------------------------------------------------------
\section{Compatibility with Existing Tokamaks}
%----------------------------------------------------------------------

The proposed mechanism requires no modification to existing tokamak
reactor designs. The deuterium gas chamber for Steps S1--S2 can be
implemented as a pre-ionization stage upstream of the plasma chamber.
The Lithium-6 blanket (Step S4) is already a planned component of
ITER and successor reactors.

The only addition is a proton beam source ($E_k \geq 1.293$ MeV)
to initiate the chain. Once initiated, Step S2 regenerates protons
internally, making the chain self-sustaining in terms of fuel.

%----------------------------------------------------------------------
\section{Open Challenges}
%----------------------------------------------------------------------

\begin{enumerate}

  \item \textbf{Chain initiation energy.}
    The initial proton beam requires $E_k \geq 1.293$ MeV.
    Standard particle accelerators provide this trivially.

  \item \textbf{Neutron thermalization.}
    The neutron produced in S1 must be thermalized before
    interacting with deuterium in S2, as the cross-section
    for $n + {}^2\text{H}$ is highest at thermal energies.
    A deuterium moderator layer achieves this naturally.

  \item \textbf{Tritium confinement.}
    Tritium produced in S2 must be confined and directed
    into the plasma chamber. This is a known engineering
    challenge already addressed in ITER design documents.

  \item \textbf{Neutrino energy loss.}
    Each S1 event loses $\sim$0.511 MeV via neutrino emission.
    This is included in the energy balance above.

\end{enumerate}

%----------------------------------------------------------------------
\section{Proposed Experimental Test}
%----------------------------------------------------------------------

\begin{enumerate}

  \item Fill a high-pressure deuterium chamber ($P > 50$ atm)
    surrounded by neutron detectors and a Lithium-6 blanket.

  \item Inject a proton beam at $E_k = 5$--20 MeV.

  \item Measure tritium production rate as a function of
    deuterium pressure. A superlinear dependence would indicate
    chain propagation beyond single-collision production.

  \item Measure heat output in the Li-6 blanket via calorimetry.

  \item If tritium production is confirmed, feed accumulated
    tritium into a small tokamak or magnetic mirror device
    and measure D--T fusion yield.

\end{enumerate}

%----------------------------------------------------------------------
\section{Conclusion}
%----------------------------------------------------------------------

We have proposed a deuterium-based neutron chain mechanism that
produces tritium \emph{in situ} via proton--electron capture,
without requiring a Coulomb barrier for the neutron production
and breeding steps. The key contributions are:

\begin{itemize}
  \item A self-sustaining proton--neutron--tritium production loop
    in dense deuterium gas.
  \item Net energy output of 23.307 MeV per complete cycle,
    compared to 17.6 MeV for standard D--T fusion alone.
  \item Full compatibility with existing tokamak infrastructure
    — no new reactor design required.
  \item Lithium-6 as a tritium recycler and yield enhancer,
    increasing the effective breeding ratio beyond unity.
\end{itemize}

If experimentally confirmed, this mechanism could resolve the
tritium supply bottleneck that currently limits the timeline
for commercial fusion energy.

%----------------------------------------------------------------------
\section*{Acknowledgments}
%----------------------------------------------------------------------

The author thanks the global fusion and nuclear physics communities
whose published work on tokamak design, tritium breeding, and
neutron physics informed this theoretical development.

%----------------------------------------------------------------------
\begin{thebibliography}{9}

\bibitem{aydogan2026a}
  E. Aydoğan,
  ``Coulomb-Barrier-Free Neutron Production via Proton--Electron
  Capture in Dense Hydrogen Gas with Lithium-6 Energy Extraction,''
  arXiv preprint, 2026.

\bibitem{iter2023}
  ITER Organization,
  ``ITER Technical Basis,''
  ITER Document Series, 2023.
  \url{https://www.iter.org}

\bibitem{east2023}
  J. Li et al.,
  ``Recent progress on EAST tokamak,''
  \textit{Nuclear Fusion}, vol. 63, 2023.

\bibitem{abdou2021}
  M. Abdou et al.,
  ``Blanket/first wall challenges and required R\&D on the pathway
  to DEMO,''
  \textit{Fusion Engineering and Design}, vol. 100, pp. 2--43, 2015.

\bibitem{pdg2022}
  R. L. Workman et al. (Particle Data Group),
  ``Review of Particle Physics,''
  \textit{Prog. Theor. Exp. Phys.}, 083C01, 2022.

\bibitem{knoll2000}
  G. F. Knoll,
  \textit{Radiation Detection and Measurement}, 3rd ed.
  Wiley, 2000.

\end{thebibliography}

\end{document}
